\documentclass[11pt]{article}
\usepackage[margin=1in]{geometry}
%%%%%%%%%%%%%%%%%%%%%%%%%%%%%%%%%%%%%%%%%%%%%%%%%%%%%%%%%%%%%%%%%%%%%%%%%%%%
\usepackage{amsmath,amssymb,amsfonts}
\usepackage{amsthm,latexsym,mathrsfs}
\usepackage{bm,bbm,dsfont}
\usepackage{color,graphicx}
\usepackage[latin1]{inputenc}
\usepackage[english]{babel}
\usepackage{algorithm}
\usepackage{algpseudocode}
\usepackage{pgfplots}
\pgfplotsset{compat=1.18}
\usepgfplotslibrary{groupplots}
% the colors of the locations, as in the main document
\definecolor{locNE}{HTML}{2A78D6}
\definecolor{locMW}{HTML}{EB6834}
\definecolor{locS}{HTML}{1BAF7A}
\definecolor{locW}{HTML}{EDA100}
\usepackage{xr-hyper}
\usepackage[colorlinks=true,linkcolor=blue,citecolor=blue,urlcolor=blue]{hyperref}
% labels of the main document, prefixed M-: build it first, so that its .aux exists
\externaldocument[M-]{math_bio_positive_matrices_clean}

\theoremstyle{plain}
\newtheorem{theorem}{Theorem}[section]
\newtheorem{lemma}[theorem]{Lemma}
\theoremstyle{remark}
\newtheorem{remark}[theorem]{Remark}

% section numbers in headings followed by a period
\makeatletter
\renewcommand\@seccntformat[1]{\csname the#1\endcsname.\quad}
\makeatother

\everymath{\displaystyle}
\numberwithin{equation}{section}
\renewcommand{\thesection}{T\arabic{section}}

\title{Three doses: a projected adaptive gradient method\thanks{Parts of this manuscript were edited with the assistance of an AI tool (Claude Opus 5.5, Anthropic). The authors are solely responsible for all proofs and statements.}}
\author{Zhilan \and John \and Yuliya \and Ludmil}
\date{}

\begin{document}

\maketitle

\begin{abstract}
Companion to ``Minimizing the spectral radius of a diagonally scaled positive
matrix''. Equation, theorem and section numbers without a prefix T refer to
that document. For the three-dose problem~\eqref{M-three-program}, which is not
convex, we use the adaptive gradient method Yogi with each step projected onto
the feasible set in the metric of its preconditioner. We show that its fixed
points, and its limits when they exist, are exactly the points satisfying the
Karush--Kuhn--Tucker conditions, and test it on $37$ instances.
\end{abstract}

\section{The problem}\label{t-problem}

We consider~\eqref{M-three-program} with $b_{k,i}>0$ for all $k=1:3$ and
$i=1:m$. By property~(\ref{M-simplex-product}) of Section~\ref{M-sec-reduced}
the feasible set is
$\mathcal{S}=\mathcal{S}_1\times\mathcal{S}_2\times\mathcal{S}_3$, with
$\mathcal{S}_k=\left(\Delta_{n-1}\right)^m$ the set of admissible
$\widehat{\bm u}_k=\left(\widehat u_{k,ij}\right)$, a product of $3m$
simplices, and the objective function is
\begin{equation}\label{t-F}
F(\widehat{\bm u}_1,\widehat{\bm u}_2,\widehat{\bm u}_3)
=\log\rho\left(A_1\operatorname{diag}\left(\bm g\left(\bm\chi\right)\right)\right),
\qquad \chi_{k,p}=\frac{b_{k,i}}{\sigma_{k,ij}}\,\widehat u_{k,ij},
\end{equation}
with $\bm g$ given by~\eqref{M-g-def}. With the eigenvectors $\bm v,\bm w>0$
of~\eqref{M-perron}, the argument that gives~\eqref{M-grad} gives
\begin{equation}\label{t-grad}
\frac{\partial F}{\partial\widehat u_{k,ij}}
=\frac{1}{\rho}\cdot\frac{\left(A_1^T\bm w\right)_pv_p}{\langle\bm w,\bm v\rangle}
\cdot\frac{\partial g_p}{\partial \chi_{k,p}}\cdot\frac{b_{k,i}}{\sigma_{k,ij}},
\qquad p=(i-1)n+j .
\end{equation}
Each evaluation of the gradient needs one Perron vector of
$A_1\operatorname{diag}(\bm g)$ and one of its transpose.

Unlike the single dose of Section~\ref{M-sec-single}, $F$ is not convex. The
dose $\chi_{2,p}$ enters both $\pi_{1,p}$, through
$\omega_v+\varepsilon_2\chi_{2,p}$, and $\pi_{2,p}$ in~\eqref{M-g-def}. For
$200\,000$ points $(\chi_{1,p},\chi_{2,p},\chi_{3,p})$ drawn log-uniformly from
$\left[10^{-5},10^{-1.5}\right]^3$, with the parameters of the model,
$\varepsilon_1=\varepsilon_2=\varepsilon_3=0.8$ and $\omega=\omega_v=1/182.5$,
the smallest value of $\chi_{2,p}^2\,\partial^2\log g_p/\partial \chi_{2,p}^2$
is $-0.154$, and $g_p$ increases with $\chi_{2,p}$ at $24.0$ per cent of the
points. The equality constraints therefore cannot be relaxed to inequalities as
in Section~\ref{M-sec-single}. Section~\ref{t-results} shows points satisfying
the Karush--Kuhn--Tucker conditions with different values of $F$; for a convex
differentiable function on a convex set every such point is a global
minimizer~\cite[Section 4.2.3]{2004Boyd_Vandenberghe}, so these values also
show that $F$ is not convex.

\section{Projected Yogi}\label{t-yogi}

We minimize $F$ in all $3N$ variables with the adaptive method Yogi
of~\cite{2018Zaheer_Yogi}, with each step projected onto $\mathcal{S}$ in the metric of its diagonal preconditioner.
The method uses the gradient~\eqref{t-grad} and no Hessian.

\subsection{Smoothness of the objective function}

For $\widehat{\bm u}\in\mathcal{S}$ all $\chi_{k,p}\ge 0$, so
$\pi_{1,p},\pi_{2,p}\ge 0$ and $\pi_{3,p}\ge 1$ in~\eqref{M-g-def}; the
numerator of $g_p$ is at least $1$ and its denominator at least $1$, and $g_p$
is a rational function of $\bm\chi$ with a positive denominator on a
neighbourhood of $\mathcal{S}$. Hence $g_p>0$ there, and
$A_1\operatorname{diag}(\bm g)$ is non-negative and irreducible with the
pattern of $A_1$. Its spectral radius is a simple eigenvalue, so it is an
analytic function of the entries near every such $\bm g$
(see~\cite[Section 6.3]{1985Horn_Johnson}), and $F$ is $C^\infty$ on a
neighbourhood of the compact set $\mathcal{S}$. In particular $\nabla F$,
given by~\eqref{t-grad}, is continuous and bounded on $\mathcal{S}$, and we
put
\begin{equation}\label{t-G}
G=\max_{\widehat{\bm u}\in\mathcal{S}}\left\|\nabla F\left(\widehat{\bm u}\right)\right\|_\infty<\infty .
\end{equation}

\subsection{Projection in a diagonal metric}

For a vector $\bm d\in\mathbb{R}^n$ with $\bm d>0$ and $\bm z\in\mathbb{R}^n$ let
\begin{equation}\label{t-proj}
\Pi^{\bm d}(\bm z)=\operatorname*{arg\,min}_{\bm y\in\Delta_{n-1}}\
\frac12\sum_{j=1}^n d_j\left(y_j-z_j\right)^2 .
\end{equation}
For $\mathcal{S}$, a product of $3m$ simplices, and a diagonal matrix $D>0$
of order $3N$, $\Pi^D_{\mathcal{S}}$ applies~\eqref{t-proj} to each factor
$(k,i)$ with the diagonal entries of $D$ of that factor. With $D=I$ this is
the Euclidean projection $\Pi_{\mathcal{S}}$.

\begin{lemma}\label{t-proj-prop}
Let $\bm d>0$ and $\bm z\in\mathbb{R}^n$.
\begin{enumerate}
\item[(a)] Problem~\eqref{t-proj} has a unique solution $\bm y$, and
$\bm y\in\Delta_{n-1}$ is this solution if and only if
\begin{equation}\label{t-vi}
\sum_{j=1}^n d_j\left(z_j-y_j\right)\left(x_j-y_j\right)\le 0
\quad\text{for all }\bm x\in\Delta_{n-1}.
\end{equation}
\item[(b)] The solution is $y_j=\max\left(z_j-\lambda/d_j,0\right)$, where
$\lambda$ is the unique root of
\begin{equation*}
\varphi(\lambda)=\sum_j\max\left(z_j-\lambda/d_j,0\right)=1 .
\end{equation*}
\item[(c)] $\Pi^{\bm d}(\bm z)$ is a continuous function of $(\bm z,\bm d)$
on $\mathbb{R}^n\times(0,\infty)^n$.
\end{enumerate}
\end{lemma}
\begin{proof}
(a) The objective function of~\eqref{t-proj} is continuous and strictly convex,
and $\Delta_{n-1}$ is compact and convex, so the minimizer exists and is
unique. For a convex differentiable function on a convex set, a point is a
minimizer if and only if the directional derivative towards every point of the
set is non-negative~\cite[Section 4.2.3]{2004Boyd_Vandenberghe}; here this
derivative is $\sum_j d_j(y_j-z_j)(x_j-y_j)$, which is~\eqref{t-vi}.

(b) The function $\varphi$ is continuous and non-increasing, it is
$0$ for $\lambda\ge\max_j d_jz_j$, and $\varphi(\lambda)\to\infty$ as
$\lambda\to-\infty$. Where $\varphi(\lambda)>0$ at least one term is positive,
and that term is strictly decreasing in $\lambda$, so $\varphi$ is strictly
decreasing on $\{\varphi>0\}$ and $\varphi(\lambda)=1$ has exactly one root.
For this $\lambda$ the vector $\bm y$ lies in $\Delta_{n-1}$. For
$\bm x\in\Delta_{n-1}$, $d_j(z_j-y_j)=\lambda$ if $y_j>0$, and
$d_j(z_j-y_j)=d_jz_j\le\lambda$ if $y_j=0$; hence
\begin{equation}
\sum_j d_j(z_j-y_j)(x_j-y_j)
=\lambda\sum_{y_j>0}(x_j-y_j)+\sum_{y_j=0}d_jz_jx_j
\le\lambda\sum_{y_j>0}(x_j-y_j)+\lambda\sum_{y_j=0}x_j
=\lambda\left(1-1\right)=0,
\end{equation}
using $x_j\ge 0$, $\sum_jx_j=1$ and $\sum_{y_j>0}y_j=1$. By (a), $\bm y$ is the
solution.

(c) The function $\varphi(\lambda;\bm z,\bm d)$ is jointly continuous. Let
$(\bm z_\ell,\bm d_\ell)\to(\bm z,\bm d)$, with roots $\lambda_\ell$ and
$\lambda$. For $\delta>0$, $\varphi(\lambda-\delta;\bm z,\bm d)>1>\varphi(\lambda+\delta;\bm z,\bm d)$
by the strict decrease in (b), so for $\ell$ large the same inequalities hold
with $(\bm z_\ell,\bm d_\ell)$, and $|\lambda_\ell-\lambda|<\delta$. Hence
$\lambda_\ell\to\lambda$, and the formula in (b) is continuous in
$(\lambda,\bm z,\bm d)$.
\end{proof}

The root in (b) is computed exactly: on each interval between consecutive
values of $d_jz_j$, $\varphi$ is affine in $\lambda$.

\begin{lemma}\label{t-fixed}
Let $D>0$ be diagonal of order $3N$, $\eta>0$, and
$\widehat{\bm u}\in\mathcal{S}$. The following are equivalent:
\begin{enumerate}
\item[(a)] $\widehat{\bm u}=\Pi^D_{\mathcal{S}}\left(\widehat{\bm u}-\eta D^{-1}\nabla F\left(\widehat{\bm u}\right)\right)$;
\item[(b)] $\nabla F\left(\widehat{\bm u}\right)^T\left(\bm x-\widehat{\bm u}\right)\ge 0$ for all $\bm x\in\mathcal{S}$;
\item[(c)] for each factor $(k,i)$ there is $\lambda_{k,i}$ with
$\partial F/\partial\widehat u_{k,ij}=\lambda_{k,i}$ where $\widehat u_{k,ij}>0$
and $\partial F/\partial\widehat u_{k,ij}\ge\lambda_{k,i}$ where $\widehat u_{k,ij}=0$,
\end{enumerate}
that is, $\widehat{\bm u}$ satisfies the Karush--Kuhn--Tucker conditions
of~\eqref{M-three-program}. In particular the condition does not depend on $D$
and $\eta$.
\end{lemma}
\begin{proof}
(a)$\Leftrightarrow$(b). Apply Lemma~\ref{t-proj-prop}(a) to each factor
with
\begin{equation*}
\bm y=\widehat{\bm u}_{k,i\cdot},\qquad
\bm z=\widehat{\bm u}_{k,i\cdot}-\eta D_{k,i}^{-1}\nabla_{k,i}F,
\end{equation*}
where the subscript $k,i$ denotes the entries of the factor. Then
$d_j(z_j-y_j)=-\eta\,\partial F/\partial\widehat u_{k,ij}$, and~\eqref{t-vi}
holding in every factor is the statement that
$\nabla_{k,i}F^T(\bm x_{k,i}-\widehat{\bm u}_{k,i\cdot})\ge 0$ for all
$\bm x_{k,i}\in\Delta_{n-1}$ and all $(k,i)$. Since $\mathcal{S}$ is the
product of the factors, this is (b).

(b)$\Leftrightarrow$(c). Condition (b) says that $\widehat{\bm u}_{k,i\cdot}$
minimizes the linear function $\bm c^T\bm x$, $\bm c=\nabla_{k,i}F$, over
$\Delta_{n-1}$ in every factor. The minimum of $\bm c^T\bm x$ over
$\Delta_{n-1}$ is $\min_jc_j$, and
$\bm c^T\bm y-\min_jc_j=\sum_jy_j\left(c_j-\min_lc_l\right)$, which vanishes for
$\bm y\in\Delta_{n-1}$ if and only if $y_j=0$ whenever $c_j>\min_lc_l$. This
is (c) with $\lambda_{k,i}=\min_jc_j$. Conversely, under (c),
$\lambda_{k,i}\le c_j$ for all $j$ with equality on the support of $\bm y$,
which has at least one element, so $\lambda_{k,i}=\min_jc_j$.

Finally, the Karush--Kuhn--Tucker conditions of~\eqref{M-three-program} in the
variables $\widehat{\bm u}$, with the constraints $\sum_j\widehat u_{k,ij}=1$
and $\widehat u_{k,ij}\ge 0$, are (c) with the multiplier $\lambda_{k,i}$ of
the equality constraint of the factor.
\end{proof}

With $D=I$ and $\eta=1$ we use
\begin{equation}\label{t-r}
r\left(\widehat{\bm u}\right)=\left\|\widehat{\bm u}-\Pi_{\mathcal{S}}\left(\widehat{\bm u}-\nabla F\left(\widehat{\bm u}\right)\right)\right\|_\infty
\end{equation}
as the measure of stationarity: $r\ge 0$ on $\mathcal{S}$, it is continuous,
and by Lemma~\ref{t-fixed} it vanishes exactly at the points satisfying
the Karush--Kuhn--Tucker conditions. It needs no threshold to decide which
entries of $\widehat{\bm u}$ are positive.

\subsection{The method}

Algorithm~\ref{t-alg-yogi} states the method. Squares, square roots, products,
quotients and $\operatorname{sign}$ of vectors are taken entrywise. Compared
with the method as stated in~\cite{2018Zaheer_Yogi}, the gradient is exact
rather than stochastic, and the step is followed by the projection; we start
the second moment $\bm s$ at the square of the first gradient, and divide the
first moment by $1-\beta_1^t$.

\begin{algorithm}[ht]
\caption{Projected Yogi}\label{t-alg-yogi}
\begin{algorithmic}[1]
\Require a start $\widehat{\bm u}^{(0)}\in\mathcal{S}$; a step $\eta>0$;
$\beta_1\in[0,1)$, $\beta_2\in(0,1)$, $\epsilon>0$; a tolerance $\mathrm{tol}$
and a maximal number of iterations $t_{\max}$
\State $\bm q\gets\nabla F\left(\widehat{\bm u}^{(0)}\right)$, $\bm \mu\gets\bm 0$, $\bm s\gets\bm q^2$
\For{$t=1,\ldots,t_{\max}$}
  \State $\bm q\gets\nabla F\left(\widehat{\bm u}^{(t-1)}\right)$
  \State $\bm\mu\gets\beta_1\bm\mu+(1-\beta_1)\bm q$
  \State $\bm s\gets\bm s-(1-\beta_2)\operatorname{sign}\left(\bm s-\bm q^2\right)\bm q^2$
  \State $D\gets\operatorname{diag}\left(\sqrt{\bm s}+\epsilon\right)$
  \State $\widehat{\bm u}^{(t)}\gets\Pi^D_{\mathcal{S}}\left(\widehat{\bm u}^{(t-1)}-\eta D^{-1}\bm\mu/\left(1-\beta_1^t\right)\right)$
  \If{$r\left(\widehat{\bm u}^{(t)}\right)<\mathrm{tol}$}
    \State \textbf{break}
  \EndIf
\EndFor
\State \Return $\widehat{\bm u}^{(t)}$, $\rho=\exp F\left(\widehat{\bm u}^{(t)}\right)$
\end{algorithmic}
\end{algorithm}

\begin{lemma}\label{t-yogi-limit}
Let $\widehat{\bm u}^{(t)}$, $\bm\mu$, $\bm s$ be the iterates of
Algorithm~\ref{t-alg-yogi} with $t_{\max}=\infty$ and $\mathrm{tol}=0$, and
$G$ as in~\eqref{t-G}.
\begin{enumerate}
\item[(a)] Every iterate lies in $\mathcal{S}$, and
$0\le s_l\le S:=\max\left(\left\|\bm s^{(0)}\right\|_\infty,(2-\beta_2)G^2\right)$
for every entry $l$ and every $t$. Hence
$\epsilon I\le D\le\left(\sqrt{S}+\epsilon\right)I$.
\item[(b)] If $\widehat{\bm u}^{(t)}\to\widehat{\bm u}^*$, then
$\widehat{\bm u}^*$ satisfies the Karush--Kuhn--Tucker conditions
of~\eqref{M-three-program}.
\end{enumerate}
\end{lemma}
\begin{proof}
(a) The iterates lie in $\mathcal{S}$ because they are values of
$\Pi^D_{\mathcal{S}}$. Fix an entry $l$, write $s$ for $s_l$ before the update
and $q^2\le G^2$ for $q_l^2$. If $s>q^2$, the update gives
$s-(1-\beta_2)q^2$, which lies in $\left(\beta_2q^2,s\right)$ if $q\ne 0$ and
equals $s$ if $q=0$. If $s<q^2$, it gives $s+(1-\beta_2)q^2\in\left[0,(2-\beta_2)q^2\right]$
when $s\ge 0$. If $s=q^2$, $s$ is unchanged. By induction from
$\bm s^{(0)}=\bm q^2\ge 0$, each entry stays non-negative, and it either does
not increase or ends at most $(2-\beta_2)G^2$; hence it never exceeds $S$.

(b) By continuity of $\nabla F$,
$\bm q^{(t)}=\nabla F\left(\widehat{\bm u}^{(t-1)}\right)\to\nabla F\left(\widehat{\bm u}^*\right)$.
The moment
\begin{equation*}
\bm\mu^{(t)}=(1-\beta_1)\sum_{\tau=1}^t\beta_1^{t-\tau}\bm q^{(\tau)}
\end{equation*}
is a weighted mean with weights summing to $1-\beta_1^t$, so
$\bm\mu^{(t)}/\left(1-\beta_1^t\right)\to\nabla F\left(\widehat{\bm u}^*\right)$.
By (a) the matrices $D^{(t)}$ lie in a compact set of positive diagonal
matrices, so a subsequence converges to some $D^*\ge\epsilon I$. Along this
subsequence, the continuity of Lemma~\ref{t-proj-prop}(c) in the update
of $\widehat{\bm u}^{(t)}$ gives
\begin{equation*}
\widehat{\bm u}^*=\Pi^{D^*}_{\mathcal{S}}\left(\widehat{\bm u}^*-\eta D^{*-1}\nabla F\left(\widehat{\bm u}^*\right)\right),
\end{equation*}
and Lemma~\ref{t-fixed} applies.
\end{proof}

Lemma~\ref{t-yogi-limit} does not assert that the iterates converge, nor
that $F$ decreases along them; the analysis of~\cite{2018Zaheer_Yogi} is for
unconstrained problems with stochastic gradients and does not cover the
projection. The method is local, and on the non-convex
problem~\eqref{M-three-program} its limit, when there is one, depends on the
start.

\section{Results}\label{t-results}

In all runs $\eta=10^{-2}$, $\beta_1=0.9$, $\beta_2=0.999$,
$\epsilon=10^{-8}$, $\mathrm{tol}=10^{-10}$ and $t_{\max}=2\cdot 10^5$; no run
reached $t_{\max}$. With $\eta=10^{-3}$ the run from the barycenter of the
weekly instance below had $r=1.3\cdot 10^{-7}$ after $5000$ iterations,
against convergence in $652$ iterations with $\eta=10^{-2}$.

\subsection{The weekly instance}

The instance is that of Section~S3 of the supplement: $b_{1i}$ the weekly
budget of the first dose, $b_{2i}=b_{3i}=b_{1i}/4$, and
$\sigma_{2,ij}=\sigma_{3,ij}=P_{ij}$, the numbers of partially protected
individuals, no budget for the second and third doses being given with the
data of~\cite{JZh}; $m=n=4$ and $N=16$. We run Algorithm~\ref{t-alg-yogi} from
$49$ starts: the barycenter $\overline{\bm u}$ of $\mathcal{S}$,
$\overline u_{k,ij}=1/n$, and the $3mn=48$ points obtained from it by
replacing one factor $(k,i)$ with $(1-\delta)\bm e_j+\delta\mathbbm{1}/n$,
$\delta=10^{-3}$. In the start $(k,i,j)$ all doses of type $k$ in location $i$
go, up to $\delta$, to the age group $j$.

From every start the method ends at $\rho=1.661826042009$, the values agreeing
to $10^{-12}$, with $r<10^{-10}$. Table~\ref{t-tab-yogi} gives the iterations
and the times. The slow runs are the $11$ starts with dose $2$ at a vertex in
locations $2$, $3$ and $4$, which took between $20546$ and $166447$
iterations; with $t_{\max}=2\cdot 10^4$ eight of them stop at $r$ between
$4\cdot 10^{-6}$ and $1.2\cdot 10^{-4}$, with $\rho$ at most
$3.3\cdot 10^{-7}$ relative above the final value. This is the slow final
phase of a first-order method, not a different limit.

\begin{table}[ht]
\centering
\small
\begin{tabular}{lcccc}
\hline
starts & barycenter & dose $1$ at a vertex & dose $2$ at a vertex & dose $3$ at a vertex\\
 & ($1$) & ($16$) & ($16$) & ($16$)\\
\hline
iterations & $652$ & $3159$--$15460$ & $12699$--$166447$ & $1118$--$3354$\\
time & $0.5$ s & $90$ s & $822$ s & $27$ s\\
\hline
\end{tabular}
\caption{Algorithm~\ref{t-alg-yogi} on the weekly instance from the $49$
starts; times are totals over the starts of the column. Every run ends at
$\rho=1.661826042009$.}\label{t-tab-yogi}
\end{table}

Figure~\ref{t-fig-conv} shows the relative excess $(\rho-\rho^*)/\rho^*$ over
$\rho^*=1.661826042009$ and the stationarity measure $r$ along three runs. The
relative excess falls below $10^{-8}$ after $174$ iterations from the
barycenter, after $340$ from the start $(3,1,4)$, and after $37290$ from the
start $(2,4,1)$, the slowest of Table~\ref{t-tab-yogi}, along which $r$
decreases slowly over most of the run.

\pgfplotsset{conv panel/.style={
  width=0.47\textwidth, height=0.36\textwidth,
  xmode=log, xlabel={iteration},
  ymode=log, ymin=1e-16, ymax=2,
  ymajorgrids, grid style={gray!25},
  axis lines*=left, title style={font=\small},
  tick label style={font=\scriptsize}, label style={font=\small},
  legend cell align=left,
  legend style={draw=none, fill=white, fill opacity=0.8, text opacity=1, font=\scriptsize},
  every axis plot/.append style={line width=1.0pt}}}
\begin{figure}[ht]
\centering
\begin{tikzpicture}
\begin{groupplot}[group style={group size=2 by 1, horizontal sep=1.4cm}, conv panel]
\nextgroupplot[title={relative excess}, ylabel={$(\rho-\rho^*)/\rho^*$}, legend pos=south west]
\addplot[color=locNE] table[x=it, y=excess]{fig_yogi_hist_bary.dat};
\addplot[color=locMW] table[x=it, y=excess]{fig_yogi_hist_k3.dat};
\addplot[color=locS]  table[x=it, y=excess]{fig_yogi_hist_k2.dat};
\legend{barycenter, {$(3,1,4)$}, {$(2,4,1)$}}
\nextgroupplot[title={stationarity measure}, ylabel={$r$}]
\addplot[color=locNE] table[x=it, y=r]{fig_yogi_hist_bary.dat};
\addplot[color=locMW] table[x=it, y=r]{fig_yogi_hist_k3.dat};
\addplot[color=locS]  table[x=it, y=r]{fig_yogi_hist_k2.dat};
\end{groupplot}
\end{tikzpicture}
\caption{Algorithm~\ref{t-alg-yogi} on the weekly instance from the barycenter
and from the starts $(k,i,j)=(3,1,4)$ and $(2,4,1)$: the relative excess over
$\rho^*=1.661826042009$ (left) and $r$ of~\eqref{t-r} (right) per iteration.
The iteration count is on a logarithmic scale, and the start, iteration $0$,
is not shown. Excesses below $10^{-16}$, rounding, are shown at
$10^{-16}$.}\label{t-fig-conv}
\end{figure}

\subsection{Thirty-six instances}

To test whether the barycenter is a reliable start we varied the data in $36$
instances: the first-dose budget $b_{1i}$ weekly or daily ($b_{1i}/7$); the
ratios
\begin{equation}
\frac{(b_{2i},b_{3i})}{b_{1i}}\in\left\{\left(\tfrac1{16},\tfrac1{16}\right),\left(\tfrac14,\tfrac14\right),(1,1),(4,4),\left(\tfrac12,\tfrac18\right),\left(\tfrac18,\tfrac12\right)\right\};
\end{equation}
and $\sigma_{2,ij}$, $\sigma_{3,ij}$ equal to $P_{ij}$, or equal to
$P_{ij}e^{Z_{ij}/2}$ and $P_{ij}e^{Z'_{ij}/2}$ with independent standard
normal $Z$, $Z'$, for two seeds. For each instance $\rho_{\mathrm{best}}$ is
the smallest value known: that of $69$ runs of the trust region interior point
method of~\cite{1999Byrd_Hribar_Nocedal} in all variables, from the $49$
starts above and $20$ random points of $\mathcal{S}$, or that of
Algorithm~\ref{t-alg-yogi} from the barycenter if it is smaller. A run is
successful if its $\rho$ is at most $(1+10^{-7})\rho_{\mathrm{best}}$.

From the barycenter, Algorithm~\ref{t-alg-yogi} takes between $545$ and
$10897$ iterations, $71$~s for all $36$ instances, and ends with $r<10^{-10}$
in every instance. It is successful in $31$ instances, $14$ of the $18$
weekly and $17$ of the $18$ daily. In two of them its value is smaller than
that of all $69$ interior point runs, by $0.13$ and $0.025$ per cent, and is
$\rho_{\mathrm{best}}$. Table~\ref{t-tab-yogi-inst} lists these two and the
five instances in which it is not successful. Four of the five are weekly,
two of them with the smallest budgets $b_{2i}=b_{3i}=b_{1i}/16$; the misses are
up to $1.2$ per cent. At these five endpoints $r<10^{-10}$, so they are points
satisfying the Karush--Kuhn--Tucker conditions to this accuracy with $\rho$
above $\rho_{\mathrm{best}}$, and $F$ is not convex on these instances.

\begin{table}[ht]
\centering
\small
\begin{tabular}{lllccc}
\hline
budget & $(b_{2i},b_{3i})/b_{1i}$ & $\sigma_2,\sigma_3$ & interior point &
Alg.~\ref{t-alg-yogi} & relative excess\\
\hline
weekly & $(\tfrac14,\tfrac14)$       & seed $2$ & $1.7117509539$ & $1.7095128048$ & $0$\\
weekly & $(\tfrac18,\tfrac12)$       & $P$      & $1.3088277638$ & $1.3084969664$ & $0$\\
\hline
weekly & $(\tfrac1{16},\tfrac1{16})$ & $P$      & $2.5031932862$ & $2.5334636776$ & $1.2\cdot 10^{-2}$\\
weekly & $(\tfrac1{16},\tfrac1{16})$ & seed $2$ & $2.4252749378$ & $2.4383599028$ & $5.4\cdot 10^{-3}$\\
weekly & $(\tfrac18,\tfrac12)$       & seed $1$ & $1.2363297040$ & $1.2376346772$ & $1.1\cdot 10^{-3}$\\
weekly & $(\tfrac18,\tfrac12)$       & seed $2$ & $1.3770469379$ & $1.3819571439$ & $3.6\cdot 10^{-3}$\\
daily  & $(\tfrac1{16},\tfrac1{16})$ & $P$      & $4.0391397841$ & $4.0402962781$ & $2.9\cdot 10^{-4}$\\
\hline
\end{tabular}
\caption{The instances in which Algorithm~\ref{t-alg-yogi} from the barycenter
gives $\rho_{\mathrm{best}}$ below the best of $69$ interior point runs (top),
and those in which it is not successful (bottom). The columns give the best
$\rho$ of the interior point runs, $\rho$ of Algorithm~\ref{t-alg-yogi}, and
its relative excess $(\rho-\rho_{\mathrm{best}})/\rho_{\mathrm{best}}$.}\label{t-tab-yogi-inst}
\end{table}

Figures~\ref{t-fig-alloc-s2} and~\ref{t-fig-alloc-p} show, for the two
instances at the top of Table~\ref{t-tab-yogi-inst}, the allocations
$\chi_{k,ij}$ at the endpoint of Algorithm~\ref{t-alg-yogi} and at that of the
interior point method from the barycenter, which gives the best of the $69$
interior point values. The two allocations differ mainly in one location. In
the first instance the interior point endpoint gives the group $60+$ of the
South $\chi_{2}=0.0162$ and $\chi_{3}=0.0211$ per day, and that of
Algorithm~\ref{t-alg-yogi} gives this group no second or third doses and their
budget to the groups $0$--$59$ of the South. In the second instance the same
holds for the group $60+$ of the West, with $\chi_{2}=0.0098$ and
$\chi_{3}=0.0181$ per day at the interior point endpoint. In the other three
locations the entries above $10^{-3}$ per day agree to $2.2$ per cent; the
smaller ones, first doses to the group $60+$, differ by up to
$3\cdot 10^{-4}$ per day. As for the single dose, where the group $60+$
receives no doses at the minimizer, the smaller value moves doses away from
this group.

\pgfplotsset{alloc panel/.style={
  width=0.46\textwidth, height=0.26\textwidth,
  xtick={0,1,2,3}, xticklabels={0--19,20--39,40--59,60+},
  ymin=0, ymajorgrids, grid style={gray!25},
  axis lines*=left, title style={font=\small},
  tick label style={font=\scriptsize}, label style={font=\small},
  legend cell align=left,
  legend style={draw=none, fill=none, font=\scriptsize},
  every axis plot/.append style={line width=1.0pt, mark size=1.8pt}}}
% #1: file stem, #2: rho of the interior point method, #3: rho of Yogi
\newcommand{\allocfigure}[3]{%
\begin{tikzpicture}
\begin{groupplot}[group style={group size=2 by 3, horizontal sep=1.2cm,
    vertical sep=1.3cm, xticklabels at=edge bottom}, alloc panel]
\nextgroupplot[title={interior point method, $\rho=#2$}, ymax=6, ylabel={$\chi_1$}, legend pos=south west]
\addplot[color=locNE, mark=*]         table[x=age, y=k1NE]{fig_alloc_#1_joint.dat};
\addplot[color=locMW, mark=square*]   table[x=age, y=k1MW]{fig_alloc_#1_joint.dat};
\addplot[color=locS,  mark=triangle*] table[x=age, y=k1S]{fig_alloc_#1_joint.dat};
\addplot[color=locW,  mark=diamond*]  table[x=age, y=k1W]{fig_alloc_#1_joint.dat};
\legend{Northeast, Midwest, South, West}
\nextgroupplot[title={Algorithm~\ref{t-alg-yogi}, $\rho=#3$}, ymax=6]
\addplot[color=locNE, mark=*]         table[x=age, y=k1NE]{fig_alloc_#1_yogi.dat};
\addplot[color=locMW, mark=square*]   table[x=age, y=k1MW]{fig_alloc_#1_yogi.dat};
\addplot[color=locS,  mark=triangle*] table[x=age, y=k1S]{fig_alloc_#1_yogi.dat};
\addplot[color=locW,  mark=diamond*]  table[x=age, y=k1W]{fig_alloc_#1_yogi.dat};
\nextgroupplot[ymax=60, ylabel={$\chi_2$}]
\addplot[color=locNE, mark=*]         table[x=age, y=k2NE]{fig_alloc_#1_joint.dat};
\addplot[color=locMW, mark=square*]   table[x=age, y=k2MW]{fig_alloc_#1_joint.dat};
\addplot[color=locS,  mark=triangle*] table[x=age, y=k2S]{fig_alloc_#1_joint.dat};
\addplot[color=locW,  mark=diamond*]  table[x=age, y=k2W]{fig_alloc_#1_joint.dat};
\nextgroupplot[ymax=60]
\addplot[color=locNE, mark=*]         table[x=age, y=k2NE]{fig_alloc_#1_yogi.dat};
\addplot[color=locMW, mark=square*]   table[x=age, y=k2MW]{fig_alloc_#1_yogi.dat};
\addplot[color=locS,  mark=triangle*] table[x=age, y=k2S]{fig_alloc_#1_yogi.dat};
\addplot[color=locW,  mark=diamond*]  table[x=age, y=k2W]{fig_alloc_#1_yogi.dat};
\nextgroupplot[ymax=100, ylabel={$\chi_3$}, xlabel={age group (years)}]
\addplot[color=locNE, mark=*]         table[x=age, y=k3NE]{fig_alloc_#1_joint.dat};
\addplot[color=locMW, mark=square*]   table[x=age, y=k3MW]{fig_alloc_#1_joint.dat};
\addplot[color=locS,  mark=triangle*] table[x=age, y=k3S]{fig_alloc_#1_joint.dat};
\addplot[color=locW,  mark=diamond*]  table[x=age, y=k3W]{fig_alloc_#1_joint.dat};
\nextgroupplot[ymax=100, xlabel={age group (years)}]
\addplot[color=locNE, mark=*]         table[x=age, y=k3NE]{fig_alloc_#1_yogi.dat};
\addplot[color=locMW, mark=square*]   table[x=age, y=k3MW]{fig_alloc_#1_yogi.dat};
\addplot[color=locS,  mark=triangle*] table[x=age, y=k3S]{fig_alloc_#1_yogi.dat};
\addplot[color=locW,  mark=diamond*]  table[x=age, y=k3W]{fig_alloc_#1_yogi.dat};
\end{groupplot}
\end{tikzpicture}}

\begin{figure}[ht]
\centering
\allocfigure{s2}{1.7117509539}{1.7095128048}
\caption{The weekly instance with $b_{2i}=b_{3i}=b_{1i}/4$ and $\sigma_2$,
$\sigma_3$ of seed $2$: the allocations $\chi_1$, $\chi_2$, $\chi_3$ (rows),
in units of $10^{-3}$ per day, at the endpoints of the interior point method
(left) and of Algorithm~\ref{t-alg-yogi} (right), both from the barycenter. One
line per location; the same scale in each row.}\label{t-fig-alloc-s2}
\end{figure}

\begin{figure}[ht]
\centering
\allocfigure{p}{1.3088277638}{1.3084969664}
\caption{As Figure~\ref{t-fig-alloc-s2}, for the weekly instance with
$b_{2i}=b_{1i}/8$, $b_{3i}=b_{1i}/2$ and $\sigma_2=\sigma_3=P$.}\label{t-fig-alloc-p}
\end{figure}

\section{Conclusions}\label{t-conclusions}

\begin{enumerate}
\item The three-dose problem~\eqref{M-three-program} is not convex: the
second dose enters $g_p$ non-convexly, and on five of the $36$ instances
Algorithm~\ref{t-alg-yogi} ends at points satisfying the
Karush--Kuhn--Tucker conditions with $\rho$ above the smallest value known.
Theorems~\ref{M-global} and~\ref{M-barrier} have no analogue here.
\item Projected Yogi, Algorithm~\ref{t-alg-yogi}, needs one gradient, that is,
two Perron vectors, per iteration and no Hessian. Its fixed points are exactly
the points satisfying the Karush--Kuhn--Tucker conditions
(Lemma~\ref{t-fixed}), and so is its limit when the iterates converge
(Lemma~\ref{t-yogi-limit}); we have no proof that they converge.
\item On the weekly instance it reaches $\rho=1.661826042009$ from all $49$
starts, with $r<10^{-10}$, in $652$ to $166447$ iterations.
\item From the barycenter it is successful in $31$ of $36$ instances, and in
two of them its value is smaller than that of $69$ interior point runs. The
misses are up to $1.2$ per cent, and the hard instances are those with small
second- and third-dose budgets.
\item In practice the method should be run from several starts, the
barycenter, the starts near the vertices and random points, and the smallest
value kept; no single start is guaranteed to give the global minimum. These
tests use second and third dose data chosen for lack of data, $\sigma_2$ and
$\sigma_3$ built from the numbers of partially protected individuals.
\end{enumerate}

The computations are in the following files:
\begin{itemize}
\item \texttt{check\_block2.py}: the sampling for the second dose in
Section~\ref{t-problem};
\item \texttt{yogi\_three\_dose.py}: Algorithm~\ref{t-alg-yogi} on the weekly
instance from the $49$ starts, and Table~\ref{t-tab-yogi};
\item \texttt{yogi\_instances\_three\_dose.py}: Algorithm~\ref{t-alg-yogi} on
the $36$ instances, and Table~\ref{t-tab-yogi-inst};
\item \texttt{test\_barycenter\_three\_dose.py}: the $69$ interior point runs
per instance;
\item \texttt{plot\_yogi\_three\_dose.py}: the data of
Figures~\ref{t-fig-conv}, \ref{t-fig-alloc-s2} and~\ref{t-fig-alloc-p}.
\end{itemize}

\bibliographystyle{plain}
\bibliography{bib_mb}

\end{document}
